\chapter{조사 준비}

\section{역할 분담}

언어조사의 효율적인 준비와 진행을 위해 조사자를 조(組)와 팀으로 나누었다. 조는 조사 준비에 필요한 각 기능을 담당하는 단위로, 조사자 6인을 2인씩 묶어 사전조사조와 장비조, 질문지조의 3개조를 구성하였다. 더불어 언어조사의 실행 단위로서 팀을 두었으며, 조사자 6인을 3인씩 묶어 `가' 팀과 `나' 팀을 구성하였다. 한 팀에 다양한 분야에 지식이 있는 사람이 고루 배치될 수 있도록 각 조에 질문지조원 1인, 사전조사조원 1인, 장비조원 1인을 배치하였고, 실제 조사에서 각각 주 질문자, 어휘·음운 중점 보조 질문자, 문법 담당 보조 질문자의 기능을 수행하였다. 팀 사이의 원활한 소통을 위하여 각 팀에 팀장 1인을 두었으며, 팀장은 팀 내 질문지조원이 맡았다. 이상의 내용을 표로 정리하면 아래와 같다.

\begin{table}[h]
  \centering
  \begin{tblr}{
        width = 0.9\textwidth,
        colspec = {Q[c,m] X[c,m] X[c,m] X[c,m] X[c,m]},
        hlines,
        vlines,
        row{1} = {font=\sffamily\bfseries},
        column{1} = {font=\sffamily\bfseries},
    }
    \diagbox{팀}{조} & 사전조사조 & 장비조 & 질문지조 & 교수·인솔자 \\
    `가' 팀 & {유승민\\(언어학과 24)} & {이채현\\(언어학과 21)} & {박준영*\\(언어학과 23)} & \SetCell[r=2]{c} {최계영(교수자)\\강민하(인솔자)\\김유빈(보조 인솔자)} \\
    `나' 팀 & {문준모\\(전기·정보공학부 25)} & {전지민\\(화학부 24)} & {이상인*\\(컴퓨터공학부 23)} &
  \end{tblr}
  \caption*{* 표시는 팀장임.}
\end{table}

조와 팀 이외에도, 언어조사 과정 전반을 보조하기 위해 총무와 사관, 일정 담당자, 보고서 담당자를 두었다. 총무는 이상인(컴공 23)으로, 재정을 관리하고 집행하는 역할을 수행하였다. 사관(史官)은 유승민(언어 24)으로, 조사 준비 과정, 조사 진행 과정, 조사 이후 회의 등을 하는 과정을 텍스트와 사진 등의 기록으로 남기는 역할을 수행하였다. 일정 담당자는 강민하(박사수료)로, 언어조사 일정을 수립하고 숙소, 식당 등을 결정하였다. 보고서 담당자는 문준모(전정 25)로, 조사 보고서의 작성과 조판을 총괄하였다.

\section{전사 훈련}

현장 조사에 나가기에 앞서, 조사자들은 1개월간 한국어 음성을 듣고 전사(轉寫)하는 것을 훈련받았다. 전사 훈련은 다음의 세 단계로 진행되었다.

9월 11일부터 9월 22일까지는 `고향에서 온 편지' 영상의 음성을 전사하였다. `고향에서 온 편지'는 1998년부터 2000년까지 SBS에서 방영되었던 예능 프로그램 `서세원의 좋은 세상 만들기'의 코너 가운데 하나로, 어르신들이 주로 자신의 자식들에게 영상 편지를 보내는 내용으로 이루어져 있다. 조사자 6인은 이들 에피소드 가운데 하나씩을 선택하여 영상의 한국어 음성을 전사하였다. 각 조사자가 선택하여 전사한 에피소드는 아래 표와 같다.

\begin{table}[h]
  \centering
  \begin{tblr}{
        width = 0.7\textwidth,
        colspec = {X[c,m] Q[c,m] X[c,m] X[c,m]},
        hlines,
        vlines,
        row{1} = {font=\sffamily\bfseries},
    }
    작업자 & 에피소드 & 지역 & 시간(분:초) \\
    문준모 & Ep.37: 반성해라! 전화 한통 드리고!! & 경기 여주 & 11:43 \\
    박준영 & Ep.15: 당신 싸가지부터 챙겨유~~ & 충남 연기 & 11:16 \\
    유승민 & Ep.6: 쏘주 한잔 걸치고 영상편지 & 충북 진천 & 12:33 \\
    이상인 & Ep.16: 선동렬도 아는 '요요' 할아버지 & 전북 순창 & 12:10 \\
    이채현 & Ep.25: 왜 너는 나를 만나서~ & 충남 금산 & 11:15 \\
    전지민 & Ep. 19: 너 돈 없냐??? & 충남 아산 & 11:31
  \end{tblr}
\end{table}

\pagebreak{}이후 9월 23일부터 10월 1일까지는 전사 자료를 보충하는 작업을 수행하였다. 일차적으로 전사한 텍스트에 오류가 없었는지 검토하고, 표준어 대역을 추가하였다. 이를 바탕으로 대상 언어에서 나타나는 어휘와 문법 표지를 표준어형과 비교하여 목록으로 정리하였다.

10월 2일부터 10월 13일까지는 교수자가 제공한 속초 지역어 음성 파일을 청취하고, 이를 전사하였다. 음성 파일의 길이는 6분 44초였다. 전사 작업이 끝나고 난 뒤에는, 수업 시간에 모여 서로의 전사문을 살펴보며 토론하는 시간을 가졌다. 질문지조는 속초 지역어 음성을 전사하는 대신 개인 언어조사를 실시하였다. 개인 언어조사는 3절에서 다룬다.

이 과정을 통해 조사자들은 일차 자료 제작의 핵심이 되는 전사 작업을 연습하였고, 이는 이후 조사 진행 및 자료 정리의 토대가 되었다.
